%% For tips on how to write a great abstract, have a look at
%%	-	https://www.cdc.gov/stdconference/2018/How-to-Write-an-Abstract_v4.pdf (presentation, start here)
%%	-	https://users.ece.cmu.edu/~koopman/essays/abstract.html
%%	-	https://search.proquest.com/docview/1417403858
%%  - 	https://www.sciencedirect.com/science/article/pii/S037837821830402X

\begin{abstract}
Generative AI has presented society with an opportunity to create a surplus of educational content in mere seconds. While knowledge is inherently beneficial and promotes intellectual and academic growth, there still is a bridge to actual learning that cannot be crossed just through using AI and its instantaneous generative abilities (Mishra \& Anthony, 2023). Learning, in the academic landscape, is comprised of obtaining information which can be in the form of notes (Rashid \& Rigas, 2006; Devlin et al., 2019). Retaining knowledge after note-taking remains a challenge, as forgetting is inevitable and overlooked by recent AI and NLP advancements. (Klemm, 2012; Naim et al., 2020). In studying and reviewing, if a person wishes to ensure their recollection of key topics is accurate and coherent, options are limited without manually referencing source material, which can be tiring and inefficient. Assistance is needed, particularly in providing instant feedback for accuracy and mastery of notes (Karpicke \& Roediger, 2008; Sachs, 1967). This research proposes Consol, a notes-to-recollection application that uses SimCSE's contrastive sentence embeddings and semantic similarity analysis to compare how similar a person's recollection is to their original notes (Gao et al., 2021; Corley et al., 2005), and uses progress tracking techniques to encourage consistency and improvement (Deterding et al., 2011; Gaston \& Cooper, 2017). It will be a web-based application that leverages semantic similarity evaluation to assess memory recall effectively. The study aims to evaluate whether semantic-based feedback enhances learning outcomes compared to traditional recall methods. Through system testing and user performance assessment, this research will analyze the system's usability, accuracy, and effectiveness in improving memory retention (Baars et al., 2022; Tran \& Nguyen, 2022). The results aim to provide a scalable solution bridging note-taking and active learning while fostering consistent knowledge consolidation (Vișcu, 2024; Stojanovic et al., 2020).
\end{abstract}
